\section{Results}
\label{sec:results}

We report a closed-loop campaign under the case-study configuration: bottoms-impurity
specification $x_{\mathrm{spec}}=0.02$, minimum projected life $\rho_{\min}=200$~h,
soft-sensor and prognostic miscoverage $\alpha_1=\alpha_2=0.10$, and similarity budget
$\varepsilon=0.05$, which set a certified floor of $0.75$. Every number below regenerates from
a single fixed seed and is frozen as committed evidence.

\subsection{The certificate's operational value}
\label{sec:res-twoarm}

Figure~\ref{fig:two_arm} contrasts the two control arms over the campaign, both scored against
the same floor. The health-blind optimiser keeps the product in specification on essentially
every cycle, with a quality coverage of $1.000$, but its unconstrained reflux pull drives the
pump hard: the projected remaining useful life clears $\rho_{\min}$ on only $0.008$ of cycles,
so the joint safety event collapses to $S_k=0.008$ (95\% Wilson interval $[0.003, 0.021]$) and
fails the floor. The conformal health-constrained arm holds $S_k=1.000$ with a Wilson lower
bound of $0.992$, meeting the floor, and its $\psi$-budget binds on every cycle ($1.00$) where
the blind arm never binds ($0.00$). Meeting a quality specification is therefore not the same
as operating safely: a loop can be in specification yet structurally unsafe over the horizon,
and the health channel together with the budget that enforces it is what carries a per-cycle
guarantee into one that survives the campaign.

\begin{figure}[t]
  \centering
  \includegraphics[width=0.82\linewidth]{fig1_two_arm.png}
  \caption{Two-arm contrast over an operating campaign, both arms scored against the same
  certified floor. The health-blind loop meets the quality spec but its over-reflux drives
  projected remaining useful life below the floor, so the joint safety event $S_k$ collapses;
  the conformal health-constrained loop holds $S_k$ at the floor.}
  \label{fig:two_arm}
\end{figure}

\subsection{Validity under measurement delay}
\label{sec:res-deadtime}

Figure~\ref{fig:deadtime} sweeps the label delay $D_a$. The joint coverage holds $S_k=1.000$
at every delay tested, $D_a\in\{0,1,2,4,8\}$ cycles, and the soft-sensor interval coverage
stays on its $0.90$ target throughout, moving only from $0.897$ at $D_a=0$ to $0.895$ at
$D_a=8$. Delay leaves validity intact and shows up only as a small movement in the coverage
rate, consistent with the $\Delta_{\mathrm{sel},k}=0$ argument of Section~\ref{sec:o1}: the
cost of stale feedback is paid in the rate, not in validity.

\begin{figure}[t]
  \centering
  \includegraphics[width=0.82\linewidth]{fig2_deadtime.png}
  \caption{Horizon coverage versus label delay $D_a$. The joint-safety coverage $S_k$ holds the
  certified floor at every delay, and the soft-sensor interval coverage stays on target;
  validity is invariant to analyser and laboratory dead time.}
  \label{fig:deadtime}
\end{figure}

\subsection{Adaptive versus static calibration}
\label{sec:res-gamma}

Figure~\ref{fig:gamma_union}(a) varies the adaptive learning rate $\gamma$: interval coverage
stays near the $0.90$ target while the half-width tightens as $\gamma$ rises, and the
one-standard-error rule selects $\gamma=0.128$ at a mean half-width of $0.0011$.
Figure~\ref{fig:gamma_union}(b) sets this against the naive alternative. A $K$-cycle Bonferroni
interval widens from a half-width of $0.0009$ at $K=1$ to $0.0013$ at $K=10$ and is already
vacuous at $K\ge 50$, where no finite calibration quantile resolves the per-cycle level
$\alpha_1/K$; adaptive conformal inference instead holds a finite half-width at target coverage
across the same horizons \citep{gibbs2021}. Horizon coverage thus requires adaptivity, and the
static union bound cannot deliver it, which is the mechanism behind the horizon guarantee of
Section~\ref{sec:horizon}.

\begin{figure}[t]
  \centering
  \includegraphics[width=0.98\linewidth]{fig3_gamma_union.png}
  \caption{(a) Coverage versus adaptivity across the adaptive-conformal learning rate. (b) The
  naive $K$-cycle Bonferroni half-width becomes vacuous at realistic horizons, whereas adaptive
  conformal inference maintains a finite half-width at target coverage.}
  \label{fig:gamma_union}
\end{figure}
